%% =============================================================================
%% chapters/03-information-architecture.tex
%% PART I: Foundations and Concepts
%% Chapter 3: Information Architecture & Standard Message Envelopes
%% =============================================================================

\chapter[Information Architecture \& Standard Message Envelopes]{Information Architecture \&\\Standard Message Envelopes}
\label{chap:information_architecture}

The Information Architecture layer defines the canonical domain models, the task state lifecycle, message transport envelopes, clinical mapping dictionaries, and the supported HL7 FHIR Release 5 (R5) resource schemas.

\begin{figure}[htbp]
    \centering
    \begin{adjustbox}{max width=\textwidth}
        %% =============================================================================
%% diagrams/fig-pragma-state-flow.tex
%% ArchiMate Pragma Task Lifecycle & 5-Phase Checkpoint State Flow
%% =============================================================================
\begin{tikzpicture}[node distance=1.2cm and 1.2cm, auto]

    % --- Background Plate ---
    \draw[archi-plate-bg] (-0.8, 0.4) rectangle (15.8, 7.2);
    \node[archi-plate-title] at (-0.6, 7.0) {Pragma Task State Machine \& 5-Phase Checkpoint Lifecycle};

    % States
    \node[archi-data-object, fill=archimotivfill, draw=archimotivdraw, minimum width=3.4cm] (st_draft)     at (1.0, 4.0) {\archibadge{Initial State}\archiname{DRAFT / REQUESTED}\\\archidesc{Gateway Ingestion}};
    \node[archi-data-object, fill=archiappfill,   draw=archiappdraw,   minimum width=3.4cm] (st_progress)  at (7.0, 4.0) {\archibadge{Active State}\archiname{IN\_PROGRESS}\\\archidesc{Pipeline Executing}};
    \node[archi-data-object, fill=architechfill, draw=architechdraw, minimum width=3.2cm] (st_completed) at (13.5, 5.5) {\archibadge{Terminal State}\archiname{COMPLETED}\\\archidesc{All Erga Success}};
    \node[archi-data-object, fill=archiimplfill, draw=archiimpldraw, minimum width=3.2cm] (st_failed)    at (13.5, 2.2) {\archibadge{Terminal State}\archiname{FAILED}\\\archidesc{Fault / Error}};

    % Checkpoints & Stages
    \node[archi-app-process, minimum width=3.2cm] (cp_ingress)  at (4.0, 5.8) {\archibadge{Checkpoint 1}\archiname{PIPELINE\_INGRESS}\\\archidesc{Conduit Ingestion}};
    \node[archi-app-process, minimum width=3.5cm] (cp_ergon)    at (7.0, 1.5) {\archibadge{Checkpoints 2 \& 3}\archiname{PRE / POST\_ERGON}\\\archidesc{Step Transitions}};
    \node[archi-app-process, minimum width=3.2cm] (cp_complete) at (10.2, 5.5) {\archibadge{Checkpoint 4}\archiname{PIPELINE\_COMPLETE}\\\archidesc{Success Audit}};
    \node[archi-app-process, minimum width=3.2cm] (cp_fail)     at (10.2, 2.2) {\archibadge{Checkpoint 5}\archiname{PIPELINE\_FAILED}\\\archidesc{Exception Context}};

    % Transitions
    \draw[archi-triggering] (st_draft) |- (cp_ingress);
    \draw[archi-triggering] (cp_ingress) -| (st_progress);

    \draw[archi-triggering] (st_progress) -- (cp_ergon);
    \draw[archi-triggering] (cp_ergon) -- (st_progress);

    \draw[archi-triggering] (st_progress) |- (cp_complete);
    \draw[archi-triggering] (cp_complete) -- (st_completed);

    \draw[archi-triggering] (st_progress) |- (cp_fail);
    \draw[archi-triggering] (cp_fail) -- (st_failed);

\end{tikzpicture}

    \end{adjustbox}
    \caption{Canonical \texttt{Pragma} State Machine and Lifecycle Progression}
    \label{fig:info_pragma_state_flow}
\end{figure}

\section{Canonical Domain Models (\texttt{Calliope})}

Harmonia avoids coupling workflow logic to proprietary or concrete wire formats by establishing strongly-typed canonical domain entities in the \texttt{Calliope} library.

\subsection{1. \texttt{Pragma} (Task State Entity)}
\texttt{Pragma} represents the single authoritative state representation for an execution task:
\begin{lstlisting}[language=Java, caption={Canonical Pragma Model Definition in Calliope}]
public class Pragma implements Serializable {
    private String pragmaId;         // Unique UUID
    private String correlationId;    // End-to-end transaction thread
    private String causationId;      // Triggering event or message ID
    private String praxisId;         // Workflow sequence definition ID
    private PragmaStatus status;     // DRAFT, REQUESTED, IN_PROGRESS, COMPLETED, FAILED
    private Integer priority;        // 1 (Lowest) - 100 (Highest)
    private Date authoredOn;
    private Date lastModified;
    private String source;           // Originating gateway / interface
    private String destination;      // Target queue or service address
    private List<ErgonPayload> input = new ArrayList<>();
    private List<ErgonPayload> output = new ArrayList<>();
    private List<PragmaCheckpoint> checkpoints = new ArrayList<>();
    private Map<String, String> metadata = new LinkedHashMap<>();
}
\end{lstlisting}

\subsection{2. \texttt{PragmaCheckpoint} (Execution Audit Snapshot)}
Captures an immutable audit milestone recorded at pipeline transitions:
\begin{lstlisting}[language=Java, caption={PragmaCheckpoint Model Definition}]
public class PragmaCheckpoint implements Serializable {
    private String checkpointId;
    private String pragmaId;
    private String ergonId;          // Active Ergon activity
    private String stageName;        // PIPELINE_INGRESS, PRE_ERGON, POST_ERGON, etc.
    private Date timestamp;
    private PragmaStatus status;
    private String statusMessage;
    private Integer stepIndex;       // 0-based sequence step
}
\end{lstlisting}

\subsection{3. \texttt{ProviderRegistryChangePragma} (Directory Change Wrapper)}
Wraps and manages asynchronous change requests submitted via FHIR Provider Registry \texttt{POST}/\texttt{PUT} endpoints:
\begin{lstlisting}[language=Java, caption={ProviderRegistryChangePragma Model Definition}]
public class ProviderRegistryChangePragma implements Serializable {
    private final Pragma pragma;

    public String getPragmaId();
    public String getCorrelationId();
    public String getOperation();           // CREATE or UPDATE
    public String getResourceType();        // Practitioner, Organization, etc.
    public String getResourceId();          // Target FHIR logical ID
    public String getRequester();           // Submitting user / principal
    public String getSourceSystem();        // Source system identifier
    public String getIfMatch();             // Expected ETag / version header
    public String getValidationOutcome();   // Validation status / diagnostics
    public String getResultingResourceVersion();
    public Resource getRequestedResource(); // Parsed FHIR R5 resource payload
}
\end{lstlisting}

\subsection{4. Canonical Provider Registry Validation Codes}
Harmonia standardizes error diagnostic reporting using canonical issue codes embedded within FHIR \texttt{OperationOutcome} responses:

\begin{table}[htbp]
\centering
\small
\begin{tabularx}{\textwidth}{l l X}
\toprule
\textbf{Error Code} & \textbf{Issue Category} & \textbf{Architectural Description \& Diagnostic Condition} \\
\midrule
\texttt{PR-VAL-001} & Structural Error & Malformed FHIR JSON, unparseable attributes, or schema violation. \\
\texttt{PR-VAL-002} & Missing Field & Required mandatory directory fields missing (e.g. Practitioner name). \\
\texttt{PR-VAL-003} & Duplicate Identifier & Duplicate deterministic provider identifier (HPI-I, HPI-O, NPI). \\
\texttt{PR-VAL-004} & Reference Not Found & Referenced target resource does not exist in Mnemosyne storage. \\
\texttt{PR-VAL-005} & Reference Inactive & Referenced target resource is marked inactive or soft-deleted. \\
\texttt{PR-VAL-006} & Concurrency Conflict & Stale update rejected due to \texttt{If-Match} / \texttt{ETag} version mismatch. \\
\texttt{PR-VAL-007} & Invalid Endpoint & Malformed technical endpoint URL, invalid payload type, or status. \\
\texttt{PR-VAL-008} & Invalid Group & Invalid group membership type or unresolved member references. \\
\texttt{PR-VAL-009} & Unsupported Resource & Submitted resource type is not supported by Provider Registry. \\
\texttt{PR-VAL-010} & Business Rule Failure & Custom business validation failure or governance policy violation. \\
\bottomrule
\end{tabularx}
\caption{Canonical Provider Registry Validation and Error Diagnostic Codes}
\label{tab:info_provider_registry_validation_codes}
\end{table}

\section{Canonical Outbound Models (\texttt{pylai-mllp-base})}

The Pylai gateway tier standardizes outbound MLLP interactions using canonical data structures:

\begin{lstlisting}[language=Java, caption={OutboundMllpRequest Model Definition}]
public class OutboundMllpRequest implements Serializable {
    private String requestId;          // Unique dispatch UUID
    private String messageControlId;   // HL7 MSH-10 control identifier
    private String destinationId;      // Target endpoint identifier (e.g. HIS_NORTH)
    private String rawMessage;         // Raw HL7 v2 pipe-delimited message string
    private String targetQueue;        // Destination Petasos queue
    private Topic topic;               // Resolved clinical topic metadata
    private Date timestamp;            // Enqueue timestamp
}
\end{lstlisting}

\section{Petasos Messaging Transport Envelopes (\texttt{petasos-api})}

Messages flowing through the messaging backbone are wrapped in a generic, opaque transport envelope:
\begin{lstlisting}[language=Java, caption={PetasosMessage Envelope Definition}]
public class PetasosMessage implements Serializable {
    private String messageId;          // Globally unique message UUID
    private String correlationId;      // Transaction correlation identifier
    private String topic;              // Logical clinical topic
    private String payloadType;        // MIME type (e.g. application/fhir+json, application/hl7-v2)
    private byte[] payload;            // Opaque binary or UTF-8 text payload
    private Map<String, String> headers = new HashMap<>();
    private long creationTimestamp;
}
\end{lstlisting}

\section{HL7 FHIR Release 5 (R5) Resource Schemas}

Harmonia maps all clinical domain entities to 15 standardized HL7 FHIR Release 5 resource schemas:
\begin{enumerate}
    \item \textbf{\texttt{Patient}}: Care recipient demographics, identifiers, and link registries.
    \item \textbf{\texttt{Encounter}}: Inpatient admissions, outpatient visits, and emergency encounters.
    \item \textbf{\texttt{Practitioner}}: Clinical practitioner records, qualifications, and credentials.
    \item \textbf{\texttt{PractitionerRole}}: Practitioner appointments, roles, and facility assignments.
    \item \textbf{\texttt{Organization}}: Health systems, hospitals, departments, and external laboratories.
    \item \textbf{\texttt{Location}}: Physical wards, operating rooms, clinics, and beds.
    \item \textbf{\texttt{HealthcareService}}: Clinical services offered at locations.
    \item \textbf{\texttt{Endpoint}}: Technical connection details (MLLP host/port, FHIR base URL).
    \item \textbf{\texttt{Group}}: Care teams and practitioner groups.
    \item \textbf{\texttt{Observation}}: Laboratory results, vital signs, and diagnostic findings.
    \item \textbf{\texttt{DiagnosticReport}}: Multi-component diagnostic summaries and radiology reports.
    \item \textbf{\texttt{DocumentReference}}: Clinical notes, discharge summaries, and pathology attachments.
    \item \textbf{\texttt{Task}}: Execution task representation mapped directly from \texttt{Pragma}.
    \item \textbf{\texttt{AuditEvent}}: Zero-PHI audit trail capturing access, creation, and modification events.
    \item \textbf{\texttt{Provenance}}: Data lineage, source system tracking, and digital signatures.
\end{enumerate}

\newpage
